\section{Conclusion}

This thesis set out to design, build and evaluate a custom video streaming pipeline for remote train operation, as a replacement for the commercial solution the project depended on. A working GStreamer based pipeline was developed and deployed across three cameras, a measurement framework was built to capture and compare end-to-end latency down to an uncertainty of around 7~ms, and the system was evaluated across a wide range of codec, transport and network configurations. The pipeline was also designed to be extensible, allowing new cameras and configurations to be added with little effort, and it served as a stable prototype throughout the development period. The objectives set out in the introduction were therefore met, and the work provides both a usable tool for the project and a foundation for the research that follows.

\subsection{Main Findings}

The first research question asked what end-to-end latency and video quality can be achieved in a GStreamer based pipeline for remote train operation, and which codec and transport configuration performs best. The results show that low latency streaming is achievable across nearly every configuration tested. An important point for reading these results is that the passthrough runs represent the real floor of what the system delivers, while the transcoded tests are staged to give full control over the stream so that individual settings can be isolated and compared against each other. The relative differences found between the transcoded settings therefore carry over onto the passthrough case as well, so whatever a given setting gains over the transcoded baseline, that same setting with passthrough would gain over the passthrough result in the same way. With that in mind, the passthrough runs produced both the lowest end-to-end latency and the best visual quality on the BRISQUE scale, with H.264 reaching 49~ms and H.265 reaching 57~ms mean, while even the transcoded baseline H.264 ran at 75~ms mean. Among the transcoded configurations, H.265 with GPU acceleration stood out, delivering 65~ms mean end-to-end while keeping the CPU below 6 percent and matching the visual quality of the much more demanding CPU based H.265. When the constant capture and render delays are added, the best configuration lands at a glass-to-glass latency of around 80~ms.

The second research question concerned how network conditions, including transport protocol and available bandwidth, affect streaming latency and reliability. Under normal conditions the difference between the two transport protocols was small but meaningful. UDP delivered around 9~ms lower mean latency than TCP, while TCP produced cleaner images, scoring around 7 points better on BRISQUE and dropping almost no packets at all. The picture changed completely once bandwidth was constrained. UDP kept a usable latency below 100~ms even while losing close to half of its packets and with it, a lot of its image quality. TCP did the opposite. Under the same constraint it fell more than 9 seconds behind, doubled the CPU load on the sender and never recovered, since its latency does not improve once it has fallen behind. Adding encryption through SRTP had almost no measurable cost in either latency or CPU, which means security can be included in the system without penalty. The reliability of the stream therefore depends heavily on knowing the available bandwidth along the route, and for unknown or varying conditions UDP is the safer choice because it degrades gracefully and restores itself rather than collapsing.

The third research question asked about the practical interplay between latency, video quality and resource usage across the different configurations, and which configuration is most suitable for remote train operation. The results show a consistent compromise between latency and quality that has to be managed rather than solved. TCP buys quality but pays for it under bandwidth pressure, and H.265 produces around 10 BRISQUE points better quality than H.264 but costs significantly more in CPU and encoding time. CPU load was found to translate directly into sender latency, and offloading the encoding to the GPU lowered both at the same time. The transit latency was shown to depend less on raw throughput and more on how regular the packet stream is, since an encoder under stress emits packets in bursts that queue unevenly. Taking all of this together, the most suitable configuration for remote train operation is passthrough with native H.265 over UDP and RTP, and H.265 with GPU acceleration over UDP and RTP where transcoding is needed. UDP should be the default, with TCP reserved for situations where the bandwidth is known to be stable and the clearest possible image is the priority. MJPEG and CPU based H.265 transcoding were excluded due to poor results, and 720p is best kept for cases where bandwidth is the limiting factor. The one question the numbers cannot answer is whether the resolution and quality are sufficient for an operator to drive safely, and that requires a human factors evaluation.

To conclude, the custom pipeline developed in this thesis is an improvement over the commercial solution it replaces. It delivers lower end-to-end latency than the previous work in the project across every configuration except MJPEG, and it does so while giving full control over each stage of the stream and bringing the measurement uncertainty down to a level where comparison against thresholds such as the 3GPP TS 22.289 100~ms target becomes meaningful. The best configurations stay at or below that target at the end-to-end level, although whether the threshold itself is the right benchmark for this specific system and its operators is a question that the technical results alone cannot settle. For the next iteration of the project, passthrough with native H.265 over UDP and RTP is the recommended starting point.

\subsection{Future Work}

The most important next step is a proper human factors evaluation. Every threshold used in this thesis comes from standards and prior research rather than from operators working with this specific system, and the question of whether a remote driver can operate safely at the latency and quality levels measured here can only be answered by involving experienced operators in the evaluation. This is also the only way to settle whether 720p is sufficient for spotting signs, obstacles and other critical objects at speed, which is something the BRISQUE metric cannot determine on its own.

On the technical side, several configurations are worth exploring further. A passthrough H.265 run at 720p would show how low the throughput can go while keeping the same quality, and a combination of H.265 with TCP under a controlled bandwidth could be of interest given how well the low throughput of H.265 fits a constrained link. An adaptive approach that matches the encoder bitrate to the available bandwidth in real time, similar to the model by Mejías et al.~\cite{mejias2024}, would address the main weakness seen in the constrained bandwidth tests, and a system that can switch transport protocol depending on the known bandwidth along the route would let the operator trade latency against quality as conditions change.

Finally, the testing itself should be expanded. The codec results come from a single run of each configuration and from a moving car rather than an actual train, so repeating the tests over multiple runs and on a real railway line would strengthen the numbers and capture the higher speeds, vibration and tower handover behaviour that a car cannot reproduce. The results also show that upgrading the camera to a higher frame rate and the monitor to a higher refresh rate would cut the constant glass-to-glass latency by almost as much as any configuration change, which makes a hardware upgrade a simple and effective improvement worth pursuing.